\textbf{Reservoir computing for chaos.} Jaeger and Haas~\cite{jaeger2004harnessing} introduced echo state networks for predicting chaotic time series. Pathak et al.~\cite{pathak2017using} train a sparse ESN with a ridge readout and run it in closed loop; their even-node squared readout is the variant we reimplement. Griffith et al.~\cite{griffith2019forecasting} show that very sparse reservoirs, even with zero spectral radius, forecast Lorenz-63 well. Haluszczynski and R\"ath~\cite{haluszczynski2019good} analyse statistically, on the Lorenz and R\"ossler systems, both the short-term prediction and the reproduction of the long-term ``climate'' by reservoirs; they report that both vary strongly among random realisations and network topologies, and that realisations with better short-term prediction also tend to reproduce the climate better. Our hypothesis H2 asks whether that agreement between the two criteria breaks when different model families, instead of realisations of one family, are ranked. Racca and Magri~\cite{racca2021robust} propose validation strategies for choosing ESN hyperparameters; we use plain validation-trajectory grid search. Platt et al.~\cite{platt2021forecasting} and Hart~\cite{hart2023attractor} relate forecasting success and faithful attractor reconstruction to generalized synchronization and conditional Lyapunov exponents; we do not test these mechanisms. Next-generation reservoir computing~\cite{gauthier2021next} replaces the random reservoir by a nonlinear vector autoregression and reports much smaller training sets; we do not run it.

\textbf{Backpropagation-trained forecasters.} Vlachas et al.~\cite{vlachas2018data} forecast high-dimensional chaotic systems with LSTMs in a reduced-order space and compare them with Gaussian processes. Their follow-up~\cite{vlachas2020backpropagation} compares gated recurrent networks trained by backpropagation through time with reservoir computing on Lorenz-96 and Kuramoto--Sivashinsky; it reports that with full-state training data the reservoir is more accurate, captures the long-term statistics better and trains faster, while with reduced-order data large reservoirs are more likely to diverge. Our setting is the full-state one, at far smaller dimension and model size. Our GRU and MLP baselines are small, one-step-trained networks optimised with Adam~\cite{kingma2014adam}.

\textbf{Benchmarks.} Gilpin~\cite{gilpin2021chaos,gilpin2023model} collects a large set of chaotic systems and compares many model classes; our study is far narrower (two systems) but varies reservoir hyperparameters, training length and the optimiser budget of the baselines under one protocol.
